\sgaStableTarget{sga6.e06.sec.2.8}2.8.
Let @@M0@@ be a perfect complex on @@M1@@. There exists an integer @@M2@@ such that @@M3@@ is a @@M4@@-perfect complex, and @@M5@@ is isomorphic to @@M6@@. Applying 2.7 to the complex @@M7@@, the sequence of distinguished triangles constructed gives, on applying the functor @@M8@@, a sequence of distinguished triangles in @@M9@@ (it is clear that @@M10@@ transforms a @@M11@@-perfect complex into a perfect complex). Moreover, we have:
@@M12@@
(note that @@M13@@ by the flatness of @@M14@@) for @@M15@@ in @@M16@@.

We therefore obtain in @@M17@@ the equality:
@@M18@@
where the @@M19@@ are the perfect complexes on @@M20@@ defined in 2.7. It follows that @@M21@@ is generated as a @@M22@@-module by the powers of @@M23@@.

To see that @@M24@@ are linearly dependent, we consider the homomorphism: @@M25@@ which associates to the class of a locally free sheaf of finite type @@M26@@ the class of the complex equal to @@M27@@ in degree @@M28@@, and zero in the other degrees. Equation (1.9.1) then gives a relation of dependence in @@M29@@.

To see that @@M30@@ are linearly independent, we can use the @@M31@@-linear morphism @@M32@@ (IV 2.11.1). We have:
@@M33@@
@@M34@@
Given a relation of the form:
@@M35@@
we multiply it by @@M36@@, and apply @@M37@@; we then find @@M38@@; we
